\chapter{总结与展望}
\label{cha:conclusion}

\section{工作总结}
\label{sec:summary}

本文围绕云原生环境下的智能运维问题，设计并实现了一个基于多智能体协作的自主演化 SRE 原型系统 OpsEvolver。与仅在静态数据集上进行根因分类的方法不同，本文强调让智能体进入真实微服务环境，通过受控故障注入、遥测采集、根因分析和经验沉淀构建完整闭环。全文的主要工作可以概括为以下几个方面。
\begin{enumerate}
    \item 在系统设计层面，本文提出了离线演化、在线诊断与反馈闭环三位一体的总体架构，并通过 ChaosExplorer、Diagnoser、Introspector 和 DeepThinker 的角色分工，将故障演练、诊断推理与知识修补统一起来。
    \item 在工程实现层面，本文实现了容器、Pod、网络、配置和磁盘 I/O 等多个层次的故障动作，接入了 Prometheus、Jaeger、日志接口和 Kubernetes 控制接口，打通了从环境交互到诊断输出的完整链路。
    \item 在知识表示层面，本文设计了知识条目、规则库和历史库三类持久化结构，使系统不仅能保存最终结论，还能保存知识形成过程，从而支持诊断复盘与失败回流。
    \item 在实验分析层面，本文基于 \texttt{history/<namespace>/dpsk.json}、\texttt{llama.json}、\texttt{qwen.json}、对应知识库文件以及 \texttt{visualize} 模块生成的统计图，围绕应用场景、可观测性边界、告警模式、知识增长、基底模型消融和代表性案例展开讨论，验证了原型系统在多类微服务场景中组织完整闭环流程的可行性。
\end{enumerate}

从研究结果看，OpsEvolver 的核心价值不只是“用 LLM 代替人工写一份 RCA 报告”，而是尝试把 LLM 的推理能力嵌入到持续演化的运维系统中，使其能够通过演练逐步形成环境特定知识。对于复杂且不断变化的云系统而言，这种“先演练、再沉淀、再反哺线上”的思路，比一次性的静态推理更符合真实运维需求。

基于 DeepSeek-V3.2 的主实验数据，OpsEvolver 已经在三个微服务系统上累计形成 959 次故障注入、703 个有效诊断周期、971 条告警、718 条知识条目和 768 条规则。其中 Astronomy Shop 的注入可观测有效率达到 100.0\%，而 Hotel Reservation 与 Social Network 中分别存在 143 个和 113 个无告警样本，说明系统已经能够同时识别“哪些故障值得学习”与“哪些故障暂时难以被观测到”。进一步的消融实验显示，Llama 4 Maverick 17B Instruct 更有利于扩大可观测覆盖，Qwen3-32B 更倾向于高密度知识抽取，而 DeepSeek-V3.2 在 Hotel Reservation 与 Social Network 上表现出更平衡的严格命中能力。虽然这些结果仍不足以支撑工业级结论，但已经足以说明：系统具备持续积累实验资产和从历史中学习的能力，论文中提出的方法论具有现实落地基础。

\section{研究不足}
\label{sec:limitations_summary}

尽管本文完成了较为完整的原型实现，但仍存在一些不可忽视的不足。
\begin{enumerate}
    \item 本文实验主要以能力验证和案例分析为主，缺少更大规模、更长时间跨度和更严格统计意义上的量化评测。
    \item 严格根因命中率仍然偏低，说明系统虽已能较稳定地产生诊断轨迹，但在故障类型命名和目标组件落点上仍有明显提升空间。
    \item 知识检索模块仍然采用较轻量的关键词匹配策略，尚未引入向量表示、图结构依赖或更精细的相似度学习方法。
    \item 多智能体协作在原型中仍偏串行化，尚未充分发挥多路径并行探索和结果投票带来的优势。
    \item 部分模块仍处于原型阶段，例如链路异常阈值设计、无告警场景下的负载增强策略、安全控制粒度以及更复杂故障组合的自动生成策略，都有待进一步打磨。
    \item 现有系统主要聚焦“诊断与知识沉淀”，尚未深入实现自动缓解和自愈执行，因此距离真正的自治 SRE 系统仍有一定距离。
\end{enumerate}

\section{未来展望}
\label{sec:future_work}

围绕上述不足，OpsEvolver 后续仍有多条值得继续推进的研究方向。
\begin{enumerate}
    \item \textbf{增强知识检索与表示能力。} 可以引入向量数据库、服务依赖图和时间序列特征，将当前轻量级关键词检索升级为更强的语义检索与因果关联建模。
    \item \textbf{强化可观测性感知的主动学习。} 针对仍存在大量无告警注入的系统，后续可联合调整工作负载、观察窗口和阈值策略，让更多真实故障被放大为可学习样本。
    \item \textbf{扩展故障组合与长链路实验。} 后续可在现有故障动作基础上支持组合故障、级联故障和跨命名空间依赖异常，更接近真实生产事故形态。
    \item \textbf{引入更严格的在线评测基准。} 可以将 ITBench、AIOpsLab 等平台进一步接入到统一实验脚本中，形成可重复执行的量化评测流程。
    \item \textbf{完善人机协作机制。} 未来可允许专家以批注、规则修订和人工审核方式参与反馈闭环，使系统在保持自治能力的同时更符合企业运维流程。
    \item \textbf{迈向自动缓解与自愈执行。} 在保证安全性的前提下，将诊断结果与回滚、扩容、重启、流量切换等缓解动作结合，探索从“会诊断”向“会处置”的演进路径。
    \item \textbf{强化生产级安全约束。} 未来应引入更细粒度的权限控制、动作审批与结果验证机制，确保智能体在真实生产环境中的可部署性。
\end{enumerate}

\section{结语}
\label{sec:final_words}

总体而言，本文完成了一个面向云原生场景的自主演化 SRE 原型系统设计与实现。尽管距离真正可大规模生产部署的自治运维系统仍有距离，但本文已经证明：通过把混沌工程、外部知识记忆、多智能体协作和反馈学习结合起来，可以把 LLM 的推理能力嵌入到持续积累、持续复盘、持续修正的运维闭环中。本文的探索为后续构建更可靠、更智能、更具环境适配能力的运维系统奠定了基础。
